\section{怎样组织好一段话}

有的中学生作文，单从一句话、一句话来看没有什么毛病，但是若把那一句话、一句话连起来看，则前言不搭后语。一篇好文章，不仅让人感到用字、遣词、造句都挺妥贴，而且要做到句与句之间丝丝入扣，联系紧密，文理通顺，合乎逻辑。当代文学大师巴金曾谈过他少年时代熟背了《古文观止》全部文章后的体会，他说：“这么多具体东西至少可以使我明白所谓‘文章’究竟是怎么一回事，它是有条有理，顺着我们的思路联下来的。”这话说得真好，因为它道出了写好文章的奥秘。是的，好文章就是一句话紧扣着一句话、一段话紧扣着—段话而写成的。

我们怎样才能把一句话又一句话组织成文理通顺、逻辑严密的一段话呢？

\textbf{1.注意句子之间所表达的意思要合乎事实，合乎事理。}

一段话总要说明一件事情或者讲清一个道理。如果句子之间所表达的内容并不能把一件事情说得明白，或者并不能把一个道理讲得清楚，那么，这段话肯定逻辑性不强，甚至有违反逻辑的地方。例如：
\newpage

(1) 李强是个聪明好学，爱说爱笑的孩子，我很喜欢他，也有信心教好他。但是究竟能不能教好他还没有把握。

(2) 我从小就爱画画，在小学五年级时还得过图画比赛二等奖，所以我一向不喜欢数理化。进入高中以来，因为我不喜欢数理化，所以对于语文很感兴趣。

（3）王林年年是“三好生”。他不仅是我们班的学习委员，而且是校篮球队代表队的主力队员之一。在这次全市中学生篮球友谊赛中，他又为我校篮球队夺得冠军立下了汗马功劳。他这种情况，使我学习的劲头更足了。

我们看，例一就说“有信心”，又说“没有把握”，前后自相矛盾。如果把“有信心”改为“愿意”、“希望”一类词语，就会平情理了。例二前一句中的“爱画画”与“不喜欢数理化”没有必然的因果关系；后一句中的“不喜欢数理化”与“对于语文很感兴趣”没有必然的因果关系；第一句的“爱画画”与第二句的“对于语文很感兴趣”也无必然的因果关系，这就叫做强加因果关系。例三中的前提“他这种情况”与结论“使我学习的劲头更足了”之间没有必然的联系，因而推理不正确。如果把“使我学习的劲头更足了”改为“使我认识到智育与体育二者并不矛盾，而是相辅相成的”，就文理通顺了。

这就是说，一段话中几个句子的关系必须要符合逻辑推理的原则，不仅要前提正确，还要看前提结论之间是不是有必然联系，得出结论的理由是不是充足，推理步骤是不是正确。只有合乎逻辑推理、意义联系紧密的一段话，其内容才可能合乎事实，合乎事理。

语法学家把这样由一组意义上联系紧密的句子构成的一段话叫做“句群”，也正是从句子与句子之间有严密的逻辑关系这个角度来提的。而“句群”的逻辑关系往往通过一些关联词语体现出来的。因此我们为了使句群联系严密，前后紧扣，要特别注意关联词语的互相呼应，用得恰当。既不要滥用关联词语，也不要缺少必要的关联词语，更不要使关联词语同句与句之间的关系不相符合。明确了这一点，对于避免由于关联词语运用不当而造成句与句之间意义不明、关系不清的情况是大有用处的。

\textbf{2.注意句子之间的意思要连贯。}

要把一段话组织好，必须要做到句子之间的意思连贯、顺畅，否则就会使读者感到文意脱节、表达不清、句子生涩了。例如：

(1) 江虹等几个女同学，在试制电动机模型的过程中遇到了不少困难。主要是因为她们能够虚心地向老师和同学请教，并且苦干加巧干，顺利地解决了许多技术上的问题。

(2) 顾医生每天来给陈民看病，陈民总是详细地把病情告诉顾医生。陈民每天打一次针，吃一些成药。不到一个月，顾医生就把陈民的病治好了。

(3) 在“大千四化”口号的鼓舞下，大家决心在三个月内生产出这种合金，填补祖国科学领域的空白。厂党委书记不泼冷水，号召全厂技术人员群策群力，请来外厂专家给予指导。钻研专业技术，攻克科学堡垒的高潮，就在这个厂里蓬蓬勃勃地掀起来了。

(4) 由于隔个小山包，我们又是悄悄地前进，所以扮作“守军”的第三组同学起初没有发现我们。直到“守军”看到了我们才匆忙抵抗，但是已经迟了。
我们看，例一由于表达的意思有缺漏，因而上下文意不能衔接，应该在“主要是”前面加上“她们终于试制成功”一类话，句子之间的意思才能连贯
\newpage

例二是讲顾医生给陈民治病这件事，但是由于几个句子不断地改变叙述角度，使读者感到头绪零乱，可以改为：“顾医生每天来给陈民看病，总是详细地了解陈民的病情。他每天给陈民打一次针，吃一些成药，不到一个月，就把陈民的病治好了。”一般说来，在一段话中，从什么角度叙述问题，应该保持相对统一，不要随意变换叙述的角度。

例三缺少必要的衔接性词语，使这一段话中的三个句子彼此失去照应，读起来生硬、不畅。如果在“广党委书记”后面的三个分句前，分别加上“不但”、“而且”、“甚至”，再在“钻研专业技术”之前加上“于是”或“这样”之类的衔接性词语，整个一段话就能够前后呼应，文意贯通了。

例四缺少必要的交代，使人搞不清既然隔个小山包怎么会被看见呢？如果把“直到‘守军’看到了我们”改成“直到我们翻过了小山包，出现在‘守军’的面前时，他们”，句意就清楚了。

总而言之，要使句子之间的意思连贯、顺畅，至少要注意三个问题：必要的意思要完整，不能有缺漏；叙述角度要统一，不能随意改变；使用必要的关联词语，不要使各个句子失去衔接和照应。

\textbf{3. 注意句子之间的层次要清楚。}

一段话的各个句子之间的关系，要让读者感到有条不紊，眉清目楚，就必须要有明晰的层次，否则将让人感到颠三倒四，思路混乱。例如：

（1）他围着一条淡紫色长绒围巾，脚穿擦得铮亮的黄皮鞋，头戴一顶“滑雪帽”。他给人的总印象真是一个讲究穿着，风度潇洒的现代青年。你看，他谈起话来滔滔不绝，走起路来轻松自如。他不仅穿着入时，而且挺有风度呢。

（2）太阳又出来了。雨过天晴，树木格外葱郁苍翠。这时风也停了，雨也住了。池子里的水清得象一面倒映蓝天的镜子。

（3）这是一幅漫画：地面有五个深浅不一的干井；画的最下面是水层；右面的右上角有个叼着烟的人，他左手拿着铁锹，右臂挎着锄子，正抬步要走；水层的上面是黑黝黝的土层；画框下面写着题目“这下面没有水，再换个地方挖”。

我们看，例一应该先说他的“穿着”，再用“不过······而且······”这个自然过渡的递进复句说到他的“厚度”，最后归结到“总印象”上去；就是谈“穿着”，也应该“头戴”、脖“围”、“脚穿”的从头到脚的次序说。这样，由分说到总说，层次就不显得杂乱无章了。

例二语言简洁，但所写的时间顺序不对。应该先说“风停”、“雨住”，再说“太阳出来”，最后说“雨过天晴”的景象，而且应该把“树木”和“池子里的水”的景象并列到一个复句中说，才能使整个一段话层次清楚。

例三语言也颇干净，但是方位的顺序说得东一榔头西一棒，不能给人以清晰的印象。应该改用由下往上说的方位顺序，由“水层”、“土层”、“干井”，地面上的人一直写到画的题目。这样写，就不会让人感到没有条理了。
应该注意：句子之间的层次关系，除了采取例一中所有析的“先分后总”（有的则是“先总后分”或者“先总后分再总”）的格式外，还可以采用几个方面逐一说明的“并列式”来表达，或者采用由小到大、由大到小、由远及近、由近及远，从上到下，从下到上，由外到内，由内到外，由表及里，由此及彼等的“推进式”来表达。“推进式”要严格遵循一定角度的顺序，颠倒了句子的顺序就将破坏文意；“并列式”也要按照人们的习惯来安排句子的顺序，如写方位的“东、南、西、北”，写颜色的“红、黄、蓝、白、黑”等等。在各种格式中，“总分”或“分总”格式是最常见的格式。这种格式中，总有“中心句”，我们应全力把它找出来，根据这个总说的“中心句”就好组织安排分说的句子了。

\textbf{4. 作修改或调整句序的练习，是培养组织好一段话的能力的好方法。}

我们在了解怎样把一个又一个句子组织成文理通顺、逻辑紧密的一段话的基本常识之后，还应进一步作一些有关的练习，使理论与实践结合起来，就可以逐渐地把掌握的知识化为灵活运用的能力。

例如，当自己或别人写了如下的几段话后，你能否辨析其毛病在哪里并加以改正呢？

（1）为了做好这道数学题，我把一切解法都想到了，还是没有解出来。当老师给我指点之后，我才恍然大悟，我怎么就是没有想出这个解法来呢。

[\ \textbf{按：}“一切解法”与“这个解法”前后自相矛盾。应将“一切解法”改为“几乎所有的解法”。]

（2）真气人，我们班的“文明班”奖旗被二班夺走了。这都怪这学期我们抓了文化科学知识的学习，因而文明礼貌、环境卫生等方面就没有抓。

[\ \textbf{按：}抓了学习并不一定就导致不抓文明礼貌等的结果。这是强加因果关系。应将“我们抓了文化科学知识的学习”改成“我们只强调抓文化科学知识的学习而忽视精神文明的建设”。]

(3) 拔河的要点，一是劲要使得齐，二是要有坚持性。二班同一班在这两点上各有所长。二班坚持性不如一班，而一班在用劲拉绳时又没有二班齐。

[\ \textbf{按：}要点有一、二两点，而后面叙述的次序却将这两点颠倒了过来，同时前面谈“各有所长”，后面却在说各有所短了，以致叙述角度不统一。应将“各有所长”后面的几句话改成“二班在用劲拉绳时比一班齐，而一班坚持性却超过了二班”。]

再如，当把下面那样结构紧密的一段话中的句子打乱了次序，你能否按照组织好一段话的基本常识把句序重新调整好呢？

(1)

①\ 它的宽大的叶子也是片片向上，几乎没有斜生的，更不用说倒垂了。

②\ 它的干通常是丈把高，象加过人工似的，一丈之内绝无旁枝。

③\ 那是力争上游的一种树，笔直的干，笔直的枝。

④\ 它所有的丫枝一律向上，而且紧紧靠拢，也象加过人工似的，成为一束，绝不旁逸斜出。

⑤\ 这是虽在北方风雪的压迫下却保持着倔强挺立的一种树。

⑥\ 它的皮光滑而有银色的晕圈，微微泛出淡青色。

[\ \textbf{按：}根据各句中加点字词的提示和有总说有分说的句意，我们可以判断这是一种“先总后分再总”的格式；分说的顺序应是：干、枝、叶、皮；分说的总起句应是第③句，由它引起分说“笔直的干，笔直的枝”；第⑤句是由外形的分说到精神本质的总说的结笔。因此，这段话的合理顺序一定是：③②④①⑥⑤。]

(2)

①\ 到了近代，利用物候知识来研究农业生产，已经发展为一门科学，就是物候学。

②\ 古代流传下来的许多农谚就包含了丰富的物候知识。

③\ 物候知识在我国起源很早。

④\ 植物的生长荣枯，动物的养育往来，如桃花开、燕子来等自然现象，称作物候。

⑤\ 物候学记录这些自然现象，从而了解随着时节推移的气候变化对动植物的影响。

[\ \textbf{按：}根据各句中加点的词的提示和对各句句意的分析，我们可以判断这是一种句句推进的格式：先说“物候”，再说“物候知识”“起源很早”；接说“古代”“物候知识”；又说“近代”“物候知识”已发展为“物候学”；最后解说“物候学”。能否把第⑤句调到最前面呢？不能。因为那样不仅前后句意无法通畅，而且使“这些自然现象”的“这些”也没有着落了。因此，这段话的合理顺序一定是：④③②①⑤。]
\newpage